% !TeX root = ../main.tex
\section{Type System and Inference}
\label{sec:type-system}

\subsection{Typing}
% !TeX root = ../main.tex
\begin{figure}[H]
\centering
\small
\begin{mathpar}
 \inferrule*[right=Refl]{ }{\tau<:\tau}
 \and
 \inferrule*[right=Trans]{\tau<:\sigma\\\sigma<:\rho}{\tau<:\rho}
 \and
 \inferrule*[right=Real-Sub]{ }{\tyR\G<:\tyR\E}
 \and
 \inferrule*[right=Arrow]{\sigma_1<:\tau_1\\\tau_2<:\sigma_2}
   {\tau_1\to\tau_2<:\sigma_1\to\sigma_2}
 \and
 \inferrule*[right=Product]{\tau_1<:\sigma_1\\\tau_2<:\sigma_2}
   {\tau_1\times\tau_2<:\sigma_1\times\sigma_2}
 \and
 \inferrule*[right=Sum]{\tau_1<:\sigma_1\\\tau_2<:\sigma_2}
   {\tau_1+\tau_2<:\sigma_1+\sigma_2}
 \and
 \inferrule*[right=List]{\tau<:\sigma}{\tyL\tau<:\tyL\sigma}
\end{mathpar}
\begin{mathpar}
 \inferrule*[right=Sub]{\typing e\tau\\\tau<:\sigma}{\typing e\sigma}
 \and
 \inferrule*[right=Real]{r\in\Real}{\typing r{\tyR m}}
 \and
 \inferrule*[right=Reject]{ }{\typing{\kw{reject}}\tau}
 \and
 \inferrule*[right=Neg]{\typing e{\tyR m}}{\typing{-e}{\tyR m}}
 \and
 \inferrule*[right=Add]{\typing{e_1}{\tyR m}\\\typing{e_2}{\tyR m}}
   {\typing{e_1+e_2}{\tyR m}}
 \and
 \inferrule*[right=Mul]{\typing{e_1}{\tyR\G}\\\typing{e_2}{\tyR m}}
   {\typing{e_1\cdot e_2}{\tyR m}}
 \and
 \inferrule*[right=Div]{\typing{e_1}{\tyR m}\\\typing{e_2}{\tyR\G}}
   {\typing{e_1/e_2}{\tyR m}}
 \and
 \inferrule*[right=Cmp]{\typing{e_1}{\tyR\G}\\\typing{e_2}{\tyR\G}}
   {\typing{e_1<e_2}{\tyB}}
\end{mathpar}
\begin{mathpar}
 \inferrule*[right=Uniform]{\typing a{\tyR m}\\\typing b{\tyR m}}
   {\typing{\draw{uniform}\alpha{a,b}}{\tyR m}}
 \and
 \inferrule*[right=Gaussian]{\typing u{\tyR m}\\\typing v{\tyR\G}}
   {\typing{\draw{gaussian}\alpha{u,v}}{\tyR m}}
 \and
 \inferrule*[right=Poisson]{\typing\lambda{\tyR m}}
   {\typing{\draw{poisson}\alpha\lambda}{\tyR m}}
 \and
 \inferrule*[right=Exponential]{\typing\lambda{\tyR\G}}
   {\typing{\draw{exponential}\alpha\lambda}{\tyR m}}
 \and
 \inferrule*[right=Bernoulli]{\typing p{\tyR m}}
   {\typing{\draw{bernoulli}\alpha p}{\tyR m}}
 \and
 \inferrule*[right=Discrete]{\typing p{\tyL{\tyR m}}}
   {\typing{\draw{discrete}\alpha p}{\tyR m}}
 \and
 \inferrule*[right=Beta]{\typing a{\tyR\G}\\\typing b{\tyR\G}}
   {\typing{\draw{beta}\alpha{a,b}}{\tyR m}}
 \and
 \inferrule*[right=Gamma]{\typing k{\tyR m}\\\typing\lambda{\tyR\G}}
   {\typing{\draw{gamma}\alpha{k,\lambda}}{\tyR m}}
\end{mathpar}
\caption{Subtyping and typing rules, with $m\in\{\G,\E\}$ and
$\alpha\in\{m,\M\}$. Rules for variables, unit and Boolean literals, functions,
recursion, application, let bindings, conditionals, pairs, projections, sum
injections and case analysis, and list constructors and case analysis are omitted
because they are standard.}
\label{fig:typing}
\Description{Subtyping permits G reals where E reals are expected. Multiplication requires a G left operand, division a G denominator, and comparison two G operands. Distribution rules share the same premises for sampling and mean evaluation.}
\end{figure}


We use the type system in \cref{fig:typing} to check which samples can be replaced
by their means. Its rules restrict how $\E$ samples may be used:
\begin{itemize}
    \item The arguments of addition are at $\E$-mode when the sum itself is.
    \item Multiplication requires a $\G$ left operand and allows an $\E$ right operand. \jules{Give examples and explain the intuition for why that's ok even though the arguments may not be independent}
    \item Division requires a $\G$ denominator but allows an $\E$ numerator.
    \item The distribution rules allow $\E$ parameters where the mean is
affine in those parameters. \jules{Give examples and explain the intuition for why that's ok}
\end{itemize}

\subsection{Type Inference}
\label{sec:type-inference}

We infer types and $\E/\G$ annotations together by generating subtype
constraints (\cref{fig:inference}). 
We choose $\E$ wherever the constraints allow it, to allow more sampling to be transformed to means.

Constraint solving follows the weak-unification approach of
\citet{traytel2011subtyping}. First, we erase affinities and unify the resulting
type shapes, rejecting constructor mismatches and infinite types (occurs check). This check ensures
that the subsequent constructor expansion terminates. We then solve the original
constraints using \cref{fig:inference-solving}, to obtain the constraints on $\E/\G$.
We then propagate modes that are forced to be $\G$ or $\E$ and set the remaining unconstrained modes to $\E$.

% !TeX root = ../main.tex
\begin{figure}[H]
\centering
\small
\newcommand{\inferconstraints}[3]{\Gamma\vdash #1\Rightarrow #2\mid #3}
\textbf{Constraint generation}\leanref{constraint-generation}
\medskip

\begin{minipage}{\linewidth}
The judgment $\inferconstraints e\tau C$ infers type $\tau$ and generates
constraints $C$. Types may contain type variables and affinity variables
$a\in\{\G,\E\}$. In the table, $\inferconstraints{e_i}{\tau_i}{C_i}$.
Each displayed $a$ is fresh; the generated constraints are $\bigcup_i C_i$
together with the listed constraints.
\end{minipage}
\medskip

\renewcommand{\arraystretch}{1.25}
\begin{tabular*}{\linewidth}{@{}l@{\extracolsep{\fill}}ll@{}}
\toprule
Expression & Inferred type & Added constraints\\
\midrule
$r$ & $\tyR a$ & $\varnothing$\\
$-e_1$ & $\tyR a$ & $\tau_1<:\tyR a$\\
$e_1+e_2$ & $\tyR a$ & $\tau_1<:\tyR a,\quad\tau_2<:\tyR a$\\
$e_1\cdot e_2$ & $\tyR a$ & $\tau_1<:\tyR\G,\quad\tau_2<:\tyR a$\\
$e_1/e_2$ & $\tyR a$ & $\tau_1<:\tyR a,\quad\tau_2<:\tyR\G$\\
$e_1<e_2$ & $\tyB$ & $\tau_1<:\tyR\G,\quad\tau_2<:\tyR\G$\\
\midrule
$\kw{uniform}(e_1,e_2)$ & $\tyR a$ & $\tau_1<:\tyR a,\quad\tau_2<:\tyR a$\\
$\kw{gaussian}(e_1,e_2)$ & $\tyR a$ & $\tau_1<:\tyR a,\quad\tau_2<:\tyR\G$\\
$\kw{discrete}(e_1)$ & $\tyR a$ & $\tau_1<:\tyL{\tyR a}$\\
\bottomrule
\end{tabular*}
\medskip

\begin{minipage}{\linewidth}
At a sampling site, $a$ also supplies the inferred annotation.
An explicit annotation $m\in\{\G,\E\}$ adds $a=m$.
\end{minipage}
\par\vspace{12pt}

\textbf{Application and let bindings}
\begin{mathpar}
\inferrule*[right=App]
  {\inferconstraints{e_1}{\tau_1}{C_1}\\
   \inferconstraints{e_2}{\tau_2}{C_2}\\
   \sigma,\rho\text{ fresh}}
  {\inferconstraints{e_1\,e_2}{\rho}
    {C_1\cup C_2\cup\{\tau_1<:(\sigma\to\rho),\ \tau_2<:\sigma\}}}
\and
\inferrule*[right=Let]
  {\inferconstraints{e_1}{\tau_1}{C_1}\\
   \Gamma,x:\tau_1\vdash e_2\Rightarrow\tau_2\mid C_2}
  {\inferconstraints{\letin x{e_1}{e_2}}{\tau_2}{C_1\cup C_2}}
\end{mathpar}
\caption{Selected constraint-generation rules. Other distributions use the
parameter types in \cref{fig:typing}.}
\label{fig:inference}
\Description{Constraint generation for arithmetic, three representative sampling
primitives, application, and let bindings. Each operand's inferred type must be
a subtype of its required type. Sampling sites share their inferred real affinity
with their annotation; explicit annotations fix that affinity.}
\end{figure}

% !TeX root = ../main.tex
\begin{figure}[H]
\centering
\small
\textbf{Subtype constraint reduction}
\medskip

\begin{minipage}{\linewidth}
After the shape check, reduce the generated constraints below.
Resolve bound type variables before each reduction and apply substitutions
throughout the constraints.
\end{minipage}
\[
\begin{aligned}
 \tyR a<:\tyR b
   &\;\leadsto\;\{a\le b\},\\[4pt]
 (\tau_1\to\tau_2)<:(\sigma_1\to\sigma_2)
   &\;\leadsto\;\{\sigma_1<:\tau_1,\;\tau_2<:\sigma_2\},\\[4pt]
 (\tau_1\star\tau_2)<:(\sigma_1\star\sigma_2)
   &\;\leadsto\;\{\tau_1<:\sigma_1,\;\tau_2<:\sigma_2\}
       &&(\star\in\{\times,+\}),\\[4pt]
 \tyL\tau<:\tyL\sigma
   &\;\leadsto\;\{\tau<:\sigma\},\\[4pt]
 \chi<:\chi
   &\;\leadsto\;\varnothing
       &&(\chi\in\{\tyU,\tyB\}).
\end{aligned}
\]
\par\vspace{8pt}

\textbf{Constructor expansion}
\medskip

\begin{minipage}{\linewidth}
$X,Y$ are unbound type variables. For a constructor-headed type $\tau$,
let $h(\tau)$ keep only its outer constructor, with fresh children:
\end{minipage}
\[
\begin{aligned}
 h(\tyR a)&=\tyR b,
 &h(\tau\to\sigma)&=X\to Y,
 &h(\tyL\tau)&=\tyL X,\\[3pt]
 h(\tau\star\sigma)&=X\star Y\quad(\star\in\{\times,+\}),
 &h(\tyU)&=\tyU,
 &h(\tyB)&=\tyB.
\end{aligned}
\]
Variables introduced by $h$ are fresh.
\smallskip

\renewcommand{\arraystretch}{1.3}
\begin{tabular*}{\linewidth}{@{}l@{\extracolsep{\fill}}ll@{}}
\toprule
Constraint & Substitution & Remaining constraint\\
\midrule
$X<:X$ & --- & none\\
$X<:Y\quad(X\ne Y)$ & --- & defer $X<:Y$\\
$X<:\tau$ & $X\mapsto h(\tau)$ & $h(\tau)<:\tau$\\
$\tau<:X$ & $X\mapsto h(\tau)$ & $\tau<:h(\tau)$\\
\bottomrule
\end{tabular*}
\medskip

\begin{minipage}{\linewidth}
The last two rows require constructor-headed $\tau$; the two occurrences of
$h(\tau)$ in a row denote the same fresh type. Revisit deferred constraints
after each expansion.
\end{minipage}
\par\vspace{12pt}

\textbf{Affinity solving}\leanref{affinity-solving}
\medskip

\begin{minipage}{\linewidth}
With $\G\le\E$, propagate assignments until no further change:
\end{minipage}
\begin{mathpar}
 \inferrule{a\le b\\a=\E}{b:=\E}
 \and
 \inferrule{a\le b\\b=\G}{a:=\G}
\end{mathpar}
\begin{minipage}{\linewidth}
Explicit annotations and required $\G$ affinities supply the initial assignments.
Conflicting assignments reject inference. Set remaining affinity variables to
$\E$ and remaining type variables to $\tyU$.
\end{minipage}

\caption{Solving the constraints generated in \cref{fig:inference}.}
\label{fig:inference-solving}
\Description{Structural subtype constraints reduce componentwise, with
contravariance for function arguments. A preliminary shape-unification check
rejects constructor clashes and infinite types. Unbound type variables expand
to constructors with fresh children; variable inequalities remain deferred. Affinity
constraints propagate E forward and G backward; remaining affinities default
to E, and remaining type variables to unit.}
\end{figure}

